% ACC 2027 / PaperPlaza official class. ieeeconf.cls MUST sit in the project
% root next to this file - Overleaf does not search subfolders for .cls.
\documentclass[letterpaper, 10 pt, conference]{ieeeconf}

\IEEEoverridecommandlockouts   % required before \thanks in ieeeconf
\overrideIEEEmargins           % required to meet the ACC printer margins

\usepackage{amsmath,amssymb,bm}
\let\proof\relax
\let\endproof\relax
\usepackage{amsthm} % for proposition/remark environments
\usepackage{booktabs,array,cite}
\usepackage{graphicx}
\usepackage{float} % enables [H] (and helps with class restrictions on lone [h])
\usepackage[hidelinks]{hyperref}
% Figures live under figures/..., so keep that on the graphics path.
\graphicspath{{figures/}{figures/disturbance/}{../experiments/nengo_nef_replication_2026-08-29/results/plots/}}

% [v6] TikZ support for the architecture figures.
\usepackage{tikz}
\usepackage{verbatim}
\usetikzlibrary{arrows.meta,positioning,fit,backgrounds,calc}
% Architecture drawing helpers live in architectures/tikz/.
\input{architectures/tikz/arch_style.tex}

\newcommand{\R}{\mathbb R}
\newcommand{\T}{^{\mathsf T}}
\newcommand{\sat}{\operatorname{sat}}
\newcommand{\clip}{\operatorname{clip}}
\newcommand{\sgnp}{\operatorname{sgn}_{+}}

% Public code/data repository (update URL for the final version).
\newcommand{\coderepo}{\url{https://github.com/your-repo}}

\input{certificate_macros.tex}
\newtheorem{proposition}{Proposition}
\newtheorem{remark}{Remark}

% [v12] Restructured around Introduction -> Problem Statement -> Online
% Neuromorphic Residual Controller -> Structural Properties and Mechanism ->
% Experimental Evaluation -> Conclusions. The local-ISS theorem and its
% numerical certificate (beta_max several orders of magnitude below the
% tested beta=0.5) are removed from this version; the theoretical centerpiece
% is now a modest bounded-residual-interface property that holds
% unconditionally at the operating point actually tested.
\title{\LARGE \bf Online PES Decoder Learning in Neuromorphic Residual UAV Attitude Control}

% ACC/PaperPlaza template author format (ieeeconf.cls does not support IEEEtran blocks).
\author{Gandhi Contreras$^{1}$, Steven Sandoval$^{1}$, and Luis Rodolfo Garcia Carrillo$^{2}$%
\thanks{$^{1}$New Mexico State University (NMSU), Las Cruces, NM, USA.}%
\thanks{$^{2}$ Space Vehicles Directorate, Air Force Research Laboratory (AFRL), Kirtland Air Force Base (KAFB), Albuquerque, NM 87108.
    Email: {\tt\small luis\_rodolfo.garcia\_carrillo@afrl.af.mil}}%
}

\usepackage{etoolbox}
\apptocmd{\thebibliography}{\itemsep=-1pt\relax}{}{}
% [v11] Layout only: tighter float separation and fuller float pages. Template fonts unchanged.
\setlength{\textfloatsep}{6pt plus 2pt minus 2pt}
\setlength{\floatsep}{5pt plus 2pt minus 2pt}
\setlength{\intextsep}{5pt plus 2pt minus 2pt}
\setlength{\dbltextfloatsep}{6pt plus 2pt minus 2pt}
\setlength{\dblfloatsep}{5pt plus 2pt minus 2pt}
\setlength{\abovecaptionskip}{3pt}
\setlength{\belowcaptionskip}{0pt}
\renewcommand{\topfraction}{0.92}
\renewcommand{\bottomfraction}{0.5}
\renewcommand{\textfraction}{0.08}
\renewcommand{\floatpagefraction}{0.8}
\renewcommand{\dbltopfraction}{0.92}
\renewcommand{\dblfloatpagefraction}{0.8}
\begin{document}
% IEEE reference practice: lists of more than six authors print as "et al."
\bstctlcite{IEEEexample:BSTcontrol}
\maketitle

% NOTE: The official ACC/PaperPlaza template typically suppresses page numbers for
% submission. Uncomment the next two lines only for internal review.
% \pagestyle{plain}
% \thispagestyle{plain}

% [v7] Rewritten: leads with the mechanism and the result; scope limits kept
% but stated once rather than after every claim.
%We present a compact neuromorphic residual controller for quaternion attitude tracking with bounded actuator commands. A fixed quaternion PD baseline is augmented by six paired coupled membrane--gate (CMG) units: continuous, non-spiking, conductance-based neurons whose gate-derived activities form a nonlinear basis. Eighteen decoder weights adapt online via prescribed error sensitivity (PES), and the bounded residual is \emph{added} to the saturated baseline so the baseline is recovered exactly at zero authority. We give bounds on internal states, command deviation, and projected decoder weights, and we isolate three commonly conflated mechanisms: intrinsic conductance modulation, decoder learning, and residual authority.
%In ten matched 45-s MuJoCo cases spanning five missions on Skydio X2 and Crazyflie~2.0 models (the latter completed with an explicit four-rotor interface), CMG--PES reduces mean quaternion-angle error by $17.3$--$60.6\%$ relative to PD while respecting actuator limits (nonzero bound occupancy only in brief Hardstep intervals). Direct and Nengo implementations agree to $2.013\times10^{-6}$ relative mean error. In the tested gain regime, most of the improvement is carried by online decoder adaptation; toggling conductance modulation is numerically indistinguishable. A benchmark against LIF, ALIF, Izhikevich, and LIFRate bases shows the same hierarchy and exposes an effort--accuracy tradeoff: families that match CMG--PES on error spend several times its residual authority to do so. Noise reruns keep the reduction range at $18.1$--$55.8\%$ with seed standard deviation below $0.031\%$, and an intermittent-disturbance study shows learned corrections reused across episodes. Results are in simulation and establish boundedness, not closed-loop stability.

% Online neuromorphic residual control can augment a conventional attitude controller under abrupt missions or unforeseen disturbances without changing its feedback gains or actuator interface. We present six continuous coupled membrane--gate (CMG) units and eighteen decoder weights adapted by prescribed error sensitivity (PES), adding a bounded correction to the saturated quaternion proportional-derivative (PD) torque request before rotor-force allocation. Across ten matched 45-second MuJoCo cases on the Skydio X2 and a Crazyflie~2 derived model, mean quaternion-angle error decreases by 17.3--60.6\%, including approximately 53--57\% reductions on the Full3 and Hard45 multi-axis tracking missions. Direct and Nengo implementations agree closely. Benchmarks against LIF, ALIF, Izhikevich, and LIFRate neuron models expose an effort-accuracy tradeoff: CMG--PES attains better tracking with substantially lower realized residual effort. Under synthetic measurement noise, the error reductions persist, and learned corrections are reused across repeated disturbances. Overall, PD provides a nominal, interpretable reference, while the online residual adds bounded, revocable adaptation that improves robustness to disturbances and noisy measurements.

\begin{comment}
    The abstract of a paper should contain a description of the problem addressed, a statement of the
main result, a key conclusion that you want to highlight, and perhaps some comment on an application/
numerical example included. One or two sentences for each of these items should be sucient.
The abstract should not include literature reviews, should not attempt to justify why the problem is
important/timely, and should not discuss problems not solved in the paper.
Abstracts are often included in searchable databases and therefore should be self-contained, with
neither bibliographic citations nor references to gures/tables in the paper. Since they are often automatically
converted to pure text, equations and non-text symbols should be avoided.
The above paragraphs are not an appropriate abstract for this paper. A more suitable abstract would
be:
\The general structure of a technical paper in the control systems area is discussed. Our aim
is to assist students/researchers that are starting to write technical papers for peer reviewed
conferences or journals. We provide examples of what to include and what not to include in
each section of a technical paper."
Note that I did not start any sentence with \In this paper, . . . " This construction is generally
redundant and one can simply drop these three words. This also applies to the introduction and remaining
sections of the paper.
\end{comment}

\begin{abstract}
% 1. Description of problem addressed.
Quadrotor attitude controllers must track time-varying references while respecting coupled rotor-force limits under modeling uncertainty, measurement noise, and disturbances.
% 2. Statement of the main result.
We present an online neuromorphic residual: six continuous membrane--gate units generate temporal features and eighteen decoder weights adapt through prescribed error sensitivity (PES). The bounded correction is added to a saturated PD quaternion controller before rotor-force allocation, preserving the classical feedback path; a bounded-residual-interface property guarantees a known per-axis command-deviation bound and exact baseline recovery at zero authority.
% 3. highlight conclusion and application/numerical examples
Against tuned PD-feedforward baselines at $\rho=3$, the residual retains geometric-mean error reductions of $6.5$--$35.6\%$, degrading only on Hardstep. The results indicate that persistent activity with online decoder learning is the principal mechanism.
\end{abstract}

%%%%%%%%%%%%%%%%%%%%%%%%%%%%%%%%%%%%
\section{Introduction}
%%%%%%%%%%%%%%%%%%%%%%%%%%%%%%%%%%%%


% 1. Should briefly describe the problem. The need to study the problem addressed.
Classical quaternion-PD feedback is an interpretable, stabilizing backbone for
tracking time-varying attitude references under coupled rotor-force limits,
but modeling uncertainty, measurement noise, and abrupt disturbances can leave
structured errors that are difficult to remove without retuning or increasing
actuator demand. Residual control addresses this limitation by adapting a
learned correction while preserving the classical feedback law and actuator
interface.

%Bounded authority makes the learning layer switchable; it does not by itself ensure stability or safety.

% 2. Relates the problem to the relevant literature.
Residual learning has been combined with classical controllers in robotics
\cite{johannink2019}, including geometric $\mathcal L_1$ quadrotor
augmentation \cite{wu2025l1quad}. Brain-inspired UAS controllers likewise
combine model-based feedback with learned terms
\cite{rubioscola2023limbic,rubioscola2021mambo}. Neuromorphic platforms support
stateful adaptive corrections through the Neural Engineering Framework and
online decoder learning \cite{eliasmith2004neuraleng,bekolay2014nengo}, commonly
using the local prescribed-error-sensitivity (PES) rule \cite{voelker2015pes}.
Self-tuning spiking PID control has been demonstrated for free-flying UAS
\cite{olivares2026selftuning}. However, continuous
membrane--gate features for actuator-aware quaternion residual control, and the
separate roles of neuron dynamics and decoder learning, remain insufficiently
characterized.

% 3. Contribution of the paper.
We study an online neuromorphic residual controller for quaternion attitude
tracking with individual rotor-force saturation. A quaternion
proportional--derivative (PD) controller supplies the baseline torque, six
continuous coupled membrane--gate (CMG, not a control-moment gyroscope) units generate temporal features, and
an eighteen-parameter decoder adapts online through PES. The bounded decoder
output is added to the saturated baseline before common rotor-force allocation;
setting the residual authority to zero therefore recovers the baseline exactly.
We compare the residual with both the original fixed-PD controller and a
tuned PD controller with desired-rate feedforward.

The contributions are:
\begin{enumerate}
    \item a fully specified CMG--PES residual architecture with paired-current
    encoding, continuous state dynamics, online decoder adaptation, bounded
    readout, explicit rotor-force allocation, and a bounded residual
    interface that guarantees exact baseline recovery at zero authority
    (Fig.~\ref{fig:arch});
    \item a controlled ablation separating conductance modulation, decoder
    learning, and residual authority, together with a matched neuron-family
    comparison that identifies persistent activity and decoder adaptation as
    the dominant mechanism; and
    \item a two-platform, five-mission evaluation against both a fixed PD
    baseline and a tuned PD-feedforward baseline, with measurement-noise and
    repeated-disturbance experiments, that isolates an explicit failure
    regime under abrupt Hardstep commands.
\end{enumerate}

We are not claiming that CMG dynamics alone create the improvement. We show that a bounded neuromorphic temporal basis, combined with online PES decoder adaptation, can augment a classical quaternion controller while preserving its actuator interface; the experiments then identify persistent activity and decoder adaptation as the dominant mechanism. Each CMG unit couples a leaky membrane to a gating state that modulates conductance, combining LIF-like integration with Hodgkin--Huxley-inspired dynamics \cite{hodgkin1952quantitative}; we use the resulting continuous (non-spiking) trajectory as a temporal feature basis, with conductance modulation secondary to decoder adaptation in the tested regime.
Against bandwidth-tuned PD-feedforward baselines, the residual retains positive geometric-mean gains on both platforms; degradation occurs on both platforms only for the Hardstep mission, with $63^\circ$--$111^\circ$ attitude changes commanded every $2~\mathrm{s}$ (Sec.~\ref{sec:ladder}).

The remainder of the paper is organized as follows. Section~\ref{sec:problem}
states the attitude-tracking problem, actuator interface, and control
objective. Section~\ref{sec:architecture} presents the CMG--PES residual
controller. Section~\ref{sec:properties} gives its structural properties
and mechanism, including the Nengo realization. Section~\ref{sec:results}
reports the baseline comparisons, ablations, neuron-family benchmark, noise
and disturbance experiments, and the Hardstep failure regime.
Section~\ref{sec:conclusion} summarizes the findings, limitations, and future work.
%
Code, model files, and run archives will be available at \coderepo.

\begin{figure*}
    \centering
    % Trim bottom whitespace in the standalone PDF so the caption sits closer.
    \includegraphics[width=\textwidth,trim=0 10mm 0 0,clip]{architectures/architecture_standalone.pdf}
    \caption{CMG--PES neuromorphic residual control architecture with online PES decoder learning.}
    \label{fig:arch}
\end{figure*}

\begin{comment}
\begin{figure*}[t]

\centering
\resizebox{\textwidth}{!}{\input{architectures/tikz/archA2p.tex}}
\caption{CMG--PES residual architecture. The light blocks are the fixed
classical path: error \eqref{eq:theta}, PD law \eqref{eq:pd}, and
allocation \eqref{eq:mixer} are never adapted. The heavy blocks encode the
clipped baseline as paired currents \eqref{eq:encoding}, evolve $2m$
membrane--gate units \eqref{eq:cmg}, and decode a bounded residual
\eqref{eq:readout} that is \emph{added} to the saturated baseline
\eqref{eq:residual}. Only $W_k$ adapts, through \eqref{eq:pes}; $\beta=0$
recovers the baseline exactly.}
\label{fig:arch}
\end{figure*}
\end{comment}

% [v6] fig:scale and fig:models were merged: both showed the same two
% aircraft. The old fig:scale call sites now point here.
\begin{figure*}[t]
\centering
% Fig. 2 (standalone PDF in figures/01_mujoco_missions/)
\includegraphics[width=\textwidth]{figures/01_mujoco_missions/fig2_drones_only.pdf}

\caption{Vehicles and simulation models. Left: CF2 hardware and MuJoCo
model. Centre: X2 hardware and MuJoCo model. Right: both geometries at
relative scale. The axis triads are body-frame axes: they are rigidly
attached to the vehicle and therefore rotate with it relative to the
inertial (world) frame (cf.\ $Q\!:\,\text{body}\to\text{inertial}$ in
Sec.~\ref{sec:problem}).}
\label{fig:models}
\end{figure*}

%%%%%%%%%%%%%%%%%%%%%%%%%%%%%%%%%%%%%%%%%%%%%%%%%%%%%%%%%%%%%%%%%%%%%%%%%%
% [v6] Nomenclature. Every symbol that appears in an equation is listed
% once here; the bracketed tag is the equation that defines it.

\begin{comment}
\begin{table*}[t]
\caption{Nomenclature. Parenthesised tags give the defining equation;
vectors are bold; $\clip_{[a,b]}$ and $\sat_{\bm\ell}$ act elementwise.}
\label{tab:nomen}\centering\scriptsize
\setlength{\tabcolsep}{4pt}
\renewcommand{\arraystretch}{0.84}
\begin{tabular}{@{}l>{\raggedright\arraybackslash}p{0.185\textwidth}l>{\raggedright\arraybackslash}p{0.185\textwidth}l>{\raggedright\arraybackslash}p{0.185\textwidth}@{}}
\toprule
\multicolumn{6}{@{}l}{\textit{Attitude, plant, and notation}}\\
$q,\;q_d$ & state / reference quaternion & $q_e$ & error quaternion, $q_d\otimes q^{-1}$ & $q_{e,0},\bm q_{e,v}$ & scalar and vector parts\\
$\bm e_q$ & shortest-path error \eqref{eq:theta} & $\theta_q$ & error angle \eqref{eq:theta} & $\bm\omega$ & body angular rate\\
$J$ & inertia tensor & $Q$ & body-to-inertial rotation & $\otimes$ & quaternion product\\
$\bm\omega_d,\bm\omega_e$ & reference and error rate \eqref{eq:errkin} & $R(\cdot)$ & rotation matrix of a quaternion & $z,z_d,\dot z$ & height, reference, rate \eqref{eq:collective}\\
$\odot,\oslash$ & elementwise $\times$, $\div$ & $\sgnp$ & sign, with $\sgnp(0)=1$ & $\sigma$ & logistic function\\
$I_m$ & identity of order $m$ & $\R$ & real numbers & $\Pi$ & Euclidean projection\\
\midrule
\multicolumn{6}{@{}l}{\textit{Fixed baseline and actuator interface}}\\
$\bm\tau_0=\bm u_0$ & baseline torque \eqref{eq:pd} & $K_q,K_\omega$ & PD gains \eqref{eq:pd} & $\bm\tau_{\max}$ & per-axis torque bound\\
$\omega_n,\zeta$ & design frequency, damping & $\bm\ell$ & per-channel command bound & $\sat_{\bm\ell}$ & saturation at $\pm\bm\ell$\\
$B$ & allocation matrix \eqref{eq:allocationmatrix} & $\bm f$ & rotor forces \eqref{eq:mixer} & $T$ & collective thrust \eqref{eq:collective}\\
$\kappa_i$ & rotor reaction coefficient & $m$ & torque channels, $m=3$ & $M$ & vehicle mass\\
$\bm f_{\max}$ & per-rotor force bound & $g$ & gravity, $9.81$\,m/s$^2$ & $\bm 1$ & vector of ones\\
\midrule
\multicolumn{6}{@{}l}{\textit{Coupled membrane--gate unit}}\\
$v_j$ & membrane state \eqref{eq:cmg} & $r_j$ & gate state \eqref{eq:cmg} & $c_j$ & shunt conductance \eqref{eq:cmg}\\
$\bar{\bm u}$ & normalized baseline \eqref{eq:encoding} & $\bm I^{\pm},\bm I$ & paired currents \eqref{eq:encoding} & $I_b,I_g$ & current offset and gain\\
$\tau,\tau_r$ & membrane, gate time constants & $v_{\rm th}$ & gate threshold & $v_{\rm rest},v_{\rm reset}$ & rest level, reversal level\\
$\gamma$ & gate self-coupling & $\alpha_h$ & logistic sharpness & $\chi$ & conductance switch, $\chi\in\{0,1\}$\\
$g_0,g_1$ & conductance gains (distinct from $g$) & $G$ & activity gain, $G=10^5$ & $h$ & integration step\\
$b_{j,k}$ & gate target \eqref{eq:gate_step} & $r^+_{j,k},v^+_{j,k}$ & updated states \eqref{eq:mem_step} & $\lambda_{j,k},d_{j,k}$ & frozen membrane coefficients\\
\midrule
\multicolumn{6}{@{}l}{\textit{Decoder, online learning, and residual authority}}\\
$\bm a_k$ & clipped activity \eqref{eq:readout} & $W_k$ & decoder \eqref{eq:readout} & $W_0$ & initial decoder, $[I_m,-I_m]$\\
$\bm s_k$ & clipped readout \eqref{eq:readout} & $\bm y^*_k$ & PES target \eqref{eq:pes} & $\bm\varepsilon_k$ & PES error \eqref{eq:pes}\\
$k_e,k^*_\omega$ & target gains \eqref{eq:pes} & $\eta$ & per-step learning rate & $\lambda_W$ & per-step decoder decay\\
$\Pi_{\mathcal W},W_{\max}$ & projection and weight bound & $\beta$ & residual authority \eqref{eq:residual} & $\bm u_k,\bm\delta_k$ & applied command, deviation \eqref{eq:residual}\\
\midrule
\multicolumn{6}{@{}l}{\textit{Reported metrics}}\\
$\bar\theta_q$ & mean quaternion-angle error & $S$ & actuator-bound occupancy [\%] & $n$ & runs per reported cell\\
\bottomrule
\end{tabular}
\end{table*}
\end{comment}

%%%%%%%%%%%%%%%%%%%%%%%%%%%%%%%%%%%%%%%%%%%%%%%%%%
\section{Problem Statement}
\label{sec:problem}
%%%%%%%%%%%%%%%%%%%%%%%%%%%%%%%%%%%%%%%%%%%%%%%%%%

\subsection{Quaternion Attitude Dynamics and Tracking Error}

Let $Q\in\mathbb S^3$ be the scalar-first Hamilton quaternion mapping body to
inertial coordinates, and let $\bm\omega$ be the body angular velocity. The
rotational dynamics are \cite{schaub2018amss}
\begin{align}
 \dot Q&=\tfrac12 Q\otimes[0,\bm\omega\T]\T,\\
 \bm J\dot{\bm\omega}&=\bm\tau_{\rm act}+\bm\tau_{\rm p}
       +\bm\tau_{\rm d}-\bm\omega\times(\bm J\bm\omega).
 \label{eq:plant}
\end{align}
where $\bm J\succ0$ is the body-frame inertia, $\bm u$ is the commanded
body torque, and $\bm\tau_{\rm act}$ is the torque realized after rotor
allocation and force clipping; hence $\bm\tau_{\rm act}=\bm u$ when the
request is feasible and generally $\bm\tau_{\rm act}\neq\bm u$ otherwise.
The terms $\bm\tau_{\rm p}$ and $\bm\tau_{\rm d}$ denote modeled passive and
externally applied disturbance torques, respectively.

Using the corresponding passive quaternion $q=Q^{-1}$ \cite{diebel2006quaternion}, define
%\cite{diebel2006quaternion}
\begin{equation}
 \dot q=-\tfrac12[0,\bm\omega\T]\T\otimes q.
 \label{eq:passive}
\end{equation}
The desired attitude $q_d$ uses the same convention. Define the attitude error
\begin{align}
 q_e&=q_d\otimes q^{-1}=[q_{e,0},\bm q_{e,v}\T]\T,\\
 \bm e_q&=\sgnp(q_{e,0})\bm q_{e,v},\qquad \sgnp(0)=1,\\
 \theta_q&=2\operatorname{atan2}(\|\bm q_{e,v}\|_2,|q_{e,0}|)
             \in[0,\pi].
 \label{eq:theta}
\end{align}
The metric $\theta_q$ is invariant to the quaternion sign. Since no continuous feedback globally stabilizes rigid-body attitude \cite{chaturvedi2011rigid}, our analysis is local (global formulations require hybrid switching \cite{mayhew2011hybrid} or geometric-algebra methods \cite{rubioscola2026ga}).
%The reported error $\theta_q$ is the shortest rotation angle, invariant under $q_e\mapsto-q_e$. Sign selection introduces a discontinuity at $q_{e,0}=0$; it does not establish robust hybrid global stability \cite{mayhew2011hybrid}.

For a desired-frame rate $\bm\omega_d$ satisfying the corresponding passive kinematics, differentiation gives
\begin{align}
 \dot q_e&=\tfrac12q_e\otimes[0,\bm\omega_e\T]\T,\\
 \bm\omega_e&=\bm\omega-R(q_e)\T\bm\omega_d,\qquad
 \dot q_{e,0}=-\tfrac12\bm q_{e,v}\T\bm\omega_e.
 \label{eq:errkin}
\end{align}
These are kinematic identities (no convergence implied). Reference Euler triples parameterize $q_d$ in this passive convention.

\subsection{Actuation and Rotor-Force Constraints}
\label{sec:actuation}
For rotor positions $\bm p_i$ in body coordinates, parallel $+z$ thrust
axes, and signed reaction coefficients $\kappa_i$, define
\begin{equation}
 B=\begin{bmatrix}
 1&1&1&1\\p_{1,y}&p_{2,y}&p_{3,y}&p_{4,y}\\
 -p_{1,x}&-p_{2,x}&-p_{3,x}&-p_{4,x}\\
 \kappa_1&\kappa_2&\kappa_3&\kappa_4
 \end{bmatrix}.
 \label{eq:allocationmatrix}
\end{equation}
The signs follow $\bm p_i\times[0,0,f_i]\T$ and the actuator's local reaction torque $\kappa_i f_i$ \cite{garciacarrillo2013quad}. Here $B$ maps rotor thrusts to total thrust and body torque; $\kappa_i$ has units of meters, not angular velocity. Given a commanded collective thrust $T$ and body-torque command $\bm u$, the corresponding inverse-allocation and clipping model is
\begin{equation}
 \begin{aligned}
 \bm f_{\rm req}&=B^{-1}[T,\bm u\T]\T,\\
 \bm f_{\rm act}&=\clip_{[\bm0,\bm f_{\max}]}(\bm f_{\rm req}),\\
 [T_{\rm act},\bm\tau_{\rm act}\T]\T&=B\bm f_{\rm act}.
 \end{aligned}
 \label{eq:mixer}
\end{equation}
Here $\bm u=\bm u_0$ for PD; the CMG path in Section~\ref{sec:architecture}
modifies that requested torque before the same allocation and clipping.
Each $f_i$ is a commanded force in newtons, not a rotor-speed command.
Exact wrench realization requires $B^{-1}[T,\bm u\T]\T\in[\bm0,\bm f_{\max}]$; for four
rotors $B$ is square and invertible, so $B^{-1}=B^{\dagger}$.
A componentwise torque box is not generally contained in this
collective-dependent feasible set. Clipped inverse allocation is used,
not an optimization-based allocator \cite{monteiro2016allocation}.

Collective thrust remains classical, with $z_d=1$ m:
\begin{equation}
 T=\clip_{[0,2.5Mg]}
 \left(\frac{M[g+6(z_d-z)-4\dot z]}{\max(R(Q)_{33},0.35)}\right).
 \label{eq:collective}
\end{equation}
Here $g$ is the gravitational acceleration. The height loop reads altitude and vertical velocity; the denominator compensates tilt.
The guard $0.35$ prevents near-singular tilt compensation as $R(Q)_{33}\to0$
(i.e., near $90^\circ$ tilt) and corresponds to limiting the effective
compensation to about $\arccos(0.35)\approx69.5^\circ$. Its upper bound does
not supersede individual rotor limits. The CMG path has no direct collective
output, but can affect height through attitude, tilt compensation, and
allocation saturation.

\subsection{Classical Attitude Controller}
Let $\bm J\succ0$ denote the body-frame inertia \cite{schaub2018amss}. If principal moments $\bm j_{\rm p}$ are given in a rotated principal-axis frame, then
\begin{equation}
 \bm J=R_{BI}\operatorname{diag}(\bm j_{\rm p})R_{BI}\T,
 \label{eq:inertiaframe}
\end{equation}
with $R_{BI}$ the principal-to-body rotation. Controller gains use the reported principal moments in their given ordering, $\bm j_{\rm tune}$ (Table~\ref{tab:models}).

The quaternion PD baseline is
\begin{align}
 \bm k_q&=2\bm j_{\rm tune}\odot\bm\omega_n^2,\qquad
 \bm k_\omega=2\zeta\bm j_{\rm tune}\odot\bm\omega_n,\\
 \bm u_0&=\sat_{\bm\ell}(-\bm k_q\odot\bm e_q-\bm k_\omega\odot\bm\omega),
 \label{eq:pd}
\end{align}
with $\bm\omega_n=(4,4,3.5)\T$ rad/s and $\zeta=0.95$. The torque box is
$\mathcal U=\prod_{i=1}^3[-\ell_i,\ell_i]$, where
$\ell_i=\max(j_{{\rm tune},i}\nu_i,10^{-6}\ \mathrm{N\,m})$ and
$\bm\nu=(12,12,10)\T$ rad/s$^2$; this controller-side box is generally not rotor-feasible.

Since \eqref{eq:pd} damps absolute rate, we also test a desired-rate feedforward baseline
\begin{equation}
 \bm u_0^{\rm ff}=\sat_{\bm\ell}\bigl(-\bm k_q\odot\bm e_q-\bm k_\omega\odot(\bm\omega-\bm\omega_d)\bigr),
 \label{eq:pdff}
\end{equation}
with $\bm\omega_d$ from central differences and optional bandwidth scaling $\bm\omega_n\mapsto\rho\bm\omega_n$ (no gyro cancellation). The same gains are used for PD and its CMG augmentation.

\subsection{Control Objective}
We minimize the shortest-path attitude error $\theta_q$ for a prescribed
reference $q_d(t)$ subject to the rotor-force limits of
Sec.~\ref{sec:actuation}. Every controller compared in Sec.~\ref{sec:results}
-- fixed PD \eqref{eq:pd}, tuned PD-feedforward \eqref{eq:pdff}, and their
CMG--PES augmentations -- shares the same state measurements,
collective-thrust loop \eqref{eq:collective}, rotor allocator
\eqref{eq:mixer}, and plant \eqref{eq:plant}; a residual controller changes
only the requested attitude torque. We report the mean error $\bar\theta_q$
over a run and the relative reduction
$100(\bar\theta_{q,\rm PD}-\bar\theta_{q,\rm CMG})/\bar\theta_{q,\rm PD}$
against each baseline.

%%%%%%%%%%%%%%%%%%%%%%%%%%%%%%%%%%%%%%%%%%%%%%%%%%%%%%%%%%%%%%%%%
\section{Online Neuromorphic Residual Controller}
\label{sec:architecture}
%%%%%%%%%%%%%%%%%%%%%%%%%%%%%%%%%%%%%%%%%%%%%%%%%%%%%%%%%%%%%%%%%
This section specifies the residual controller in state-space form: a
temporal feature basis (CMG units), an online decoder update (PES), and the
bounded control law that adds the decoded correction to the classical
baseline. Section~\ref{sec:properties} then gives its structural properties.

\subsection{CMG Temporal Features}

For $m$ torque channels and $2m$ units, normalize the clipped baseline and
form positive/negative current pairs:
\begin{equation}
 \begin{aligned}
 \bar{\bm u}&=\bm u_0\oslash\bm\ell,\\
 \bm I^{\pm}&=I_b\bm1+I_g\max(\pm\bar{\bm u},\bm0),\\
 \bm I&=[(\bm I^+)\T,(\bm I^-)\T]\T.
 \end{aligned}
 \label{eq:encoding}
\end{equation}
Each unit has a dimensionless membrane state $v_j$ and gate $r_j$:
\begin{align}
 \dot r_j&=[\sigma(\alpha_h(v_j-v_{\rm th}+\gamma r_j))-r_j]/\tau_r,\\
 \dot v_j&=-(v_j-v_{\rm rest})/\tau+I_j/\tau-c_j(v_j-v_{\rm reset}),\\
 c_j&=(g_0+\chi g_1|I_j|)r_j^2,\qquad \chi\in\{0,1\}.
 \label{eq:cmg}
\end{align}
% [v6] CMG unit coupling topology (TikZ source under architectures/tikz/).
\begin{figure}[t]
\centering
\resizebox{0.86\columnwidth}{!}{\input{architectures/tikz/archB1c.tex}}
\caption{Coupling topology of one CMG unit, cf.~\eqref{eq:cmg}. The
membrane sets the gate target; the gate acts back twice, shifting its own
threshold through $\gamma$ and shunting the membrane through
$c_j\propto r_j^2$ ($\chi=1$ makes that shunt input dependent). The
dynamics are continuous: no spike, no reset, and only $a_j$ leaves the
unit.}
\label{fig:cmgunit}
\end{figure}
%
The logistic gate target depends on $v_j$, while $r_j$ shifts that target
and modulates membrane conductance (Fig.~\ref{fig:cmgunit}). Thus, both states participate even
though the decoder reads the gate-derived activity. The model has no
threshold-triggered reset or spike event; $v_{\rm reset}$ is a reversal
level in a continuous conductance term. The gains $g_0,g_1$ have the units
needed to make $c_j$ a rate.

The actual integration is sequential, not exponential Euler for both
states. From the old states compute
\begin{align}
 b_{j,k}&=\sigma\!\left(\clip_{[-60,60]}
       [\alpha_h(v_{j,k}-v_{\rm th}+\gamma r_{j,k})]\right),\\
 r_{j,k}^{+}&=\clip_{[0,1]}
       [r_{j,k}+(h/\tau_r)(b_{j,k}-r_{j,k})].
 \label{eq:gate_step}
\end{align}
Then evaluate $c_{j,k}=(g_0+\chi g_1|I_{j,k}|)(r_{j,k}^{+})^2$,
$\lambda_{j,k}=\tau^{-1}+c_{j,k}$, and
$d_{j,k}=(v_{\rm rest}+I_{j,k})/\tau+c_{j,k}v_{\rm reset}$. The
frozen-coefficient membrane step is
\begin{equation}
 v_{j,k}^{+}=e^{-\lambda_{j,k}h}v_{j,k}
       +(1-e^{-\lambda_{j,k}h})d_{j,k}/\lambda_{j,k}.
 \label{eq:mem_step}
\end{equation}
This exactly solves the scalar membrane equation with those frozen
coefficients; it does not establish unconditional stability of the coupled
nonlinear closed loop. Assign $(v_{j,k+1},r_{j,k+1})=(v_{j,k}^{+},r_{j,k}^{+})$.
The exponential-Euler step $r_{j,k}^{+}=e^{-h/\tau_r}r_{j,k}+(1-e^{-h/\tau_r})b_{j,k}$ is
convex for every $h>0$ without a clip, and changes only the weight on $r_{j,k}$
($0.905$ vs.\ $0.9$ at $h/\tau_r=0.1$); all results use \eqref{eq:gate_step}.

\subsection{PES Decoder Adaptation}
The within-step activity, initial decoder, and clipped readout are
\begin{equation}
 \begin{aligned}
 \bm a_k&=\clip_{[0,1]}(G\bm r_k^+),\qquad W_0=[I_m,-I_m],\\
 \bm s_k&=\clip_{[-1,1]}(W_k\bm a_k).
 \end{aligned}
 \label{eq:readout}
\end{equation}
Pair subtraction retains directional information from a nonnegative
activity basis; state memory does not imply instantaneous sign agreement
after a command reversal. With $G=10^5$ the activity saturates once
$r_j\ge10^{-5}$, so $G$ is a primary design hyperparameter for how much of
$r_j$'s dynamic range reaches the decoder. The teaching target and local update are
\begin{align}
 \bm y_k^*&=\clip_{[-1,1]}(-k_e\bm e_{q,k}-k_\omega^*\bm\omega_k),\\
 \bm\varepsilon_k&=\bm y_k^*-\bm s_k,\\
 W_{k+1}&=\Pi_{\mathcal W}
       [W_k+\eta\bm\varepsilon_k\bm a_k\T-\lambda_W W_k].
 \label{eq:pes}
\end{align}
Here $\mathcal W=[-W_{\max},W_{\max}]^{m\times2m}$. The coefficients $\eta$ and $\lambda_W$ are per-step (so changing $h$ changes the effective learning dynamics). PES is a local error-driven update toward the prescribed target $\bm y_k^*$; it does not differentiate through the plant or claim to minimize a closed-loop cost.

\subsection{Bounded Residual Control Law}
Using $W_k$ before the update, the requested command is
\begin{equation}
 \bm u_k=\sat_{\bm\ell}(\bm u_{0,k}+\beta\bm\ell\odot\bm s_k),\quad
 \bm\delta_k=\bm u_k-\bm u_{0,k},\quad 0\le\beta\le1.
 \label{eq:residual}
\end{equation}
The canonical UAV case has $m=3$, six CMG units, and eighteen decoder
weights. Intrinsic conductance modulation ($\chi$), synaptic/decoder
adaptation ($W_k$), and output authority ($\beta\bm\ell$) are distinct.

%%%%%%%%%%%%%%%%%%%%%%%%%%%%%%%%%%%%%%%%%%%%%%%%%%%%%%%%%%%%%%%%%
\section{Structural Properties and Mechanism}
\label{sec:properties}
%%%%%%%%%%%%%%%%%%%%%%%%%%%%%%%%%%%%%%%%%%%%%%%%%%%%%%%%%%%%%%%%%
We state only the properties that hold at the tested operating point,
including $\beta=0.5$: internal-state bounds, a decoder-weight bound from
projection, and a uniform command-deviation bound set by residual
authority. We then interpret the PES update's contraction behavior and
report the Nengo realization.

\subsection{Bounded States, Weights, and Residual Authority}
\begin{proposition}[Bounded internal states]\label{prop:states}
With the nonnegative conductance parameters in Table~\ref{tab:parameters},
$r_{j,0}\in[0,1]$, and $0<h\le\tau_r$, the gate remains in $[0,1]$.
The membrane iterate remains in the convex hull of its initial value,
$v_{\rm reset}$, and the interval
$[v_{\rm rest}+I_b,v_{\rm rest}+I_b+I_g]$.
\end{proposition}
\begin{proof}
The unclipped gate step is a convex combination of $r_{j,k}$ and
$b_{j,k}\in[0,1]$. Since $I_j\in[I_b,I_b+I_g]$ and $c_j\ge0$,
$d_j/\lambda_j$ is a convex combination of $v_{\rm rest}+I_j$ and
$v_{\rm reset}$. Equation~\eqref{eq:mem_step} is a convex combination of
$v_{j,k}$ and $d_j/\lambda_j$. Induction proves both statements.
\end{proof}

\begin{proposition}[Bounded residual interface]\label{prop:interface}
If $W_0\in\mathcal W$, then \eqref{eq:pes} keeps $W_k\in\mathcal W$ for
every $k$. For $\bm u_{0,k}\in\mathcal U$ and any $W_k\in\mathcal W$
(in particular any finite $W_k$), the applied command satisfies
$\bm u_k\in\mathcal U$ and
\begin{equation}
 |\delta_{k,i}|\le\beta\ell_i,\qquad i=1,\dots,m,
 \label{eq:interfacebound}
\end{equation}
for every $\beta\in[0,1]$, including the tested $\beta=0.5$. If $\beta=0$,
the requested baseline command is recovered exactly at every step; if
$W_k=0$ at a particular step, it is recovered at that step, but learning
need not keep $W_{k+1}=0$.
\end{proposition}
\begin{proof}
Projection gives $W_k\in\mathcal W$ for every $k$ by induction on
\eqref{eq:pes}. The outer clip in \eqref{eq:residual} enforces
$\bm u_k\in\mathcal U$; its componentwise nonexpansiveness and
$\sat_{\ell_i}(u_{0,k,i})=u_{0,k,i}$ give
$|\delta_{k,i}|\le|\beta\ell_i s_{k,i}|\le\beta\ell_i$. The recovery
statements follow by substitution.
\end{proof}
This is the theoretical centerpiece we rely on: it applies to the
architecture as actually run, holds unconditionally in $\beta$ (not only
for a conservative sub-range), and does not claim more than is proved. It
does not establish state safety, rotor headroom, closed-loop stability, or
a bound on tracking improvement; a bounded adversarial correction can
still degrade or destabilize a nonlinear plant. Framed against standard
robustness analysis, the uniform command-deviation bound
\eqref{eq:interfacebound} is the kind of bounded-perturbation premise an
input-to-state-stability argument would take as given, not a substitute
for one; we stop at that premise and do not claim the ISS conclusion a
full argument would still need to establish.

\subsection{Baseline Recovery}
Proposition~\ref{prop:interface} makes $\beta=0$ an exact, single-parameter
shutoff: the same allocator and plant state then give the same actuator
output as the fixed-PD controller, with no separate code path or decoder
reset required. The residual is therefore a strictly optional, revocable
augmentation of the classical loop rather than a replacement for it.

\subsection{Interpretation of PES Adaptation}
\begin{proposition}[Decoder contraction under fixed features]\label{prop:pes}
If the target and a nonzero activity $\bm a$ are fixed, leak is zero, and
readout clipping and projection remain inactive, then
\begin{equation}
 \bm\varepsilon_{k+1}
       =(1-\eta\|\bm a\|_2^2)\bm\varepsilon_k.
 \label{eq:contraction}
\end{equation}
For a nonzero initial error, its norm contracts to zero when
$0<\eta\|\bm a\|_2^2<2$.
\end{proposition}
\begin{proof}
Under the stated conditions,
$(W_{k+1}-W_k)\bm a=\eta\bm\varepsilon_k\|\bm a\|_2^2$;
subtracting this from the fixed target error gives
\eqref{eq:contraction}. The scalar multiplier has modulus less than one
under the stated condition.
\end{proof}
This is the constant-feature PES result \cite{voelker2015pes}, not a
closed-loop convergence theorem: changing targets, loss of excitation,
clipping, and leakage invalidate a direct application to the experiment,
and no unique identification of all decoder weights follows from one fixed
regressor. Since $\|\bm a\|^2\le2m$, $0<\eta<1/m$ is sufficient for that
restricted model with nonzero activity. Its practical content is that
contraction needs \emph{persistent}, nonzero activity: a unit whose
activity is silent ($\bm a\approx\bm0$) stalls the update regardless of
$\eta$. This is exactly the property that separates the neuron families in
Sec.~\ref{sec:family} and, in its leaky, clipped form
\eqref{eq:constantbasis}, explains the Hardstep lag of Sec.~\ref{sec:ladder}.

\subsection{Nengo Realization}
In Nengo~4.1.0 a custom \texttt{NeuronType} exposes membrane and gate states, and direct current injection bypasses NEF gain/intercept fitting; unfiltered CMG activity uses $\eta_N=\eta(2m)/h$ and error $\bm s-\bm y^*$, with leak on pre-step weights. Its teaching graph adds a one-step latency, so the weight sequences differ; over ten direct/Nengo cases the maximum relative mean-error difference is $2.013\times10^{-6}$: tested numerical agreement, not formal equivalence.

\section{Experimental Evaluation}
\label{sec:results}
\subsection{Platforms, Missions, and Metrics}

We evaluate in MuJoCo rigid-body simulation with site-applied actuator wrenches
\cite{todorov2012mujoco}, using two Menagerie-derived multirotors (CF2, X2)
\cite{menagerieModels}. The attitude loop commands body torque and the height loop
collective thrust. For CF2 we expose four rotor-force channels and use the clipped
inverse allocator \eqref{eq:mixer}. Table~\ref{tab:models} lists simulation parameters and
spec-context values; the CF2 $1.5\times$ scale increases a model limit, not measured
hardware capability. Both models omit motor/RPM/battery states and validated rotor
aerodynamics; feedback uses \texttt{qpos}/\texttt{qvel}, with noise only in
Sec.~\ref{sec:noise}.

\begin{table}[t]
\centering
\caption{Model/interface parameters (top) and canonical CMG--PES settings (bottom). Neural quantities are normalized model variables; learning coefficients are per step. Mass matches published specs \cite{bitcrazeCF20Specs,skydioX2ESpecs}; other limits/parameters are simulation values.}
\label{tab:models}\label{tab:parameters}
\setlength{\tabcolsep}{3pt}
\footnotesize
\begin{tabular}{@{}p{0.40\linewidth}p{0.27\linewidth}p{0.25\linewidth}@{}}
\toprule
Parameter & CF2 rotor model & X2 model\\
\midrule
Mass $M$ (kg) & $0.027$ & $1.325$\\
Rotor $|p_x|,|p_y|$ (m) & $0.0325,0.0325$ & $0.14,0.18$\\
Reaction magnitude $|\kappa_i|$ (m) & $0.005$ & $0.0201$\\
Rotor force limit (N) & $.0875\to.13125$ & $13$\\
Total force limit (N) & $.350\to.525$ & $52$\\
Main-test thrust/weight & $1.982$ & $4.001$\\
$\bm j_{\rm tune}$ ($10^{-5}$ kg\,m$^2$) & $(2.3951,\allowbreak2.3951,\allowbreak3.2347)$ & ---\\
$\bm j_{\rm tune}$ ($10^{-2}$ kg\,m$^2$) & --- & $(6.0711,\allowbreak3.6468,\allowbreak2.5412)$\\
Rigid-body integrator & RK4 & Euler\\
Experiment timestep (s) & $0.0005$ & $0.0005$\\
\bottomrule
\end{tabular}
\par\medskip
\footnotesize\setlength{\tabcolsep}{3pt}
\begin{tabular}{@{}ll@{\hspace{6pt}}ll@{}}
\toprule
CMG--PES & Value & & Value\\\midrule
$h$, duration & $0.0005$ s, $45$ s & $\alpha_h,\gamma,g_0,g_1$ & $50,2,120,30$\\
$\tau,\tau_r$ & $0.05$ s, $0.005$ s & $I_b,I_g,G$ & $0.75,1.25,10^5$\\
$v_{\rm rest},v_{\rm reset},v_{\rm th}$ & $0,0,1$ & $\beta,\eta$ & $0.5,8\times10^{-4}$\\
$v_{j,0},r_{j,0},\chi$ & $0,0,1$ & $\lambda_W,W_{\max},k_e,k^*_\omega$ & $10^{-6},6,8,0$\\
\bottomrule
\end{tabular}
\end{table}



CF2 body axes are principal. For X2 the geometry-derived inertia is
$\bm J=10^{-2}[3.66517,0,-.21000;\allowbreak\,0,2.54117,0;\allowbreak\,-.21000,0,6.05280]$\,kg\,m$^2$;
torque limits use $\bm j_{\rm tune}$.

Each PD/CMG pair shares initial state, reference, duration, allocator, model, and timestep. The primary campaign is deterministic (no injected disturbances or noise), one run per cell, learning active throughout.

Reference Euler triples (rad): Roll $(0.25\sin0.7t,0,0)$; Full3 $0.5\sin(t+(0,\pi/2,\pi)\T)$; Hard45
$(0.55\sin0.75t+0.20\sin(2.10t+0.40),\allowbreak\;0.45\sin(0.95t+\pi/3)+0.18\sin1.75t,\allowbreak\;0.65\sin(0.55t+\pi)+0.25\sin(1.35t+0.20))$;
Hardstep: $0$ for 1 s then 2-s cycles through $(.75,-.45,.55)$, $(-.65,.55,-.70)$, $(.50,.50,.90)$, $(-.80,-.35,.20)$;
Yaw360 $(0,0,2\pi(3\xi^2-2\xi^3))$, $\xi=t/45$. Hover is omitted (near-zero error).

\subsection{Comparison with Classical Baselines}
\label{sec:matrix}
Because \eqref{eq:pd} damps absolute rate, moving references induce a
velocity lag under fixed PD; the desired-rate feedforward baseline
\eqref{eq:pdff} separates this effect from learning, and sweeping
$\bm\omega_n\mapsto\rho\bm\omega_n$ over $\rho\in[0.75,8]$ gives a best
mission-mean at $\rho=3$ on both airframes (CF2 becomes rotor-limited
beyond $\rho\approx4$). Against this tuned PD-feedforward baseline -- the
strongest classical comparison we test -- CMG--PES retains geometric-mean
quaternion-angle error reductions of $6.5\%$ (CF2) and $35.6\%$ (X2), with
Hardstep the sole degrading case (Sec.~\ref{sec:ladder}); CMG--PES improves
each rung (PD, feedforward PD, tuned feedforward PD), but the margin
shrinks and realized residual effort drops as the baseline strengthens.
The CF2/X2 asymmetry at $\rho=3$ (geometric-mean reductions of $6.5\%$
versus $35.6\%$) is consistent with this authority effect: CF2's
rotor-limited onset at $\rho\approx4$ means its $\rho=3$ feedforward
baseline is already close to consuming the actuator headroom the residual
would otherwise use, whereas X2 retains more margin at the same $\rho$.
Read this way, the shrinking benefit is where classical actuator authority
is already saturated, not a platform-specific failure of the basis; we
have not isolated this mechanism from a direct headroom measurement, so we
state it as the reading consistent with the data rather than as
demonstrated.

Against the original fixed PD law \eqref{eq:pd}, the wider comparison used
throughout the rest of this section, Table~\ref{tab:all} shows CMG--PES
lowering mean quaternion-angle error in all ten nominal cases by
$17.3$--$60.6\%$ while sharing the actuator interface; the largest
reductions occur on Yaw360 and Full3 and the smallest on Roll. Holding the
decoder at $W_k=W_0$ recovers at most a fifth of that gain
(Sec.~\ref{sec:ablation}). Table~\ref{tab:all} also reports four other
neuron families under the same baseline; Sec.~\ref{sec:family} summarizes
that comparison.

\input{tables/table_all}

\begin{figure}[t]
\centering
\includegraphics[width=\linewidth]{figures/02_neural_activity/attitude_detail/x2_hard45_attitude.png}
\caption{X2 Hard45: PD baseline versus CMG--PES over the 45-s run. Hard45 drives
all three axes at once and gives a large reduction that is also insensitive to
measurement noise (Table~\ref{tab:all}).}
\label{fig:tracking}
\end{figure}

% Internal-signal figure removed to meet the 8-page ACC limit.
% (Move to supplementary material / repository if needed.)
%\begin{figure}[t]
%\centering
%\includegraphics[width=0.98\columnwidth]{figures/09_CMGfixed_vs_CMGPES/x2_hard45_fixed_vs_pes.png}
%\caption{X2 Hard45: internal CMG signals for Fig.~\ref{fig:tracking} (fixed decoder vs. PES).}
%\label{fig:full3_neurons}
%\end{figure}

Figure~\ref{fig:tracking} illustrates typical tracking on X2 Hard45.
Improvement refers to the mean metric; it need not hold at every instant.

%%%%%%%%%%%%%%%%%%%%%%%%%%%%%%%%%%%%%%%
\subsection{Ablation: What Drives the Improvement?}\label{sec:ablation}

We ablate decoder adaptation (PES), conductance modulation ($\chi$), and readout gain ($G$). With decoder mode held fixed, toggling $\chi$ (retaining the base shunt $g_0r_j^2$) is numerically indistinguishable.

\emph{Fixed versus adaptive decoder.} Holding $W_k=W_0$ keeps CMG dynamics but limits the readout to the paired difference $\bm a^+-\bm a^-$. It yields only modest reductions (Table~\ref{tab:all}) compared to PES.

\emph{Gate/activity regime.} At $I_j=I_b$, $v_j\approx0.75<v_{\rm th}$ and $r_j\approx\sigma(-12.5)\approx3.7\times10^{-6}$, i.e., a small but nonzero logistic tail. Once $r_j\ge1/2$, Proposition~\ref{prop:states} implies $b_{j,k}\ge1/2$ in \eqref{eq:gate_step}, so opened gates latch; with $G=10^5$ the resting activity is $a_j\approx0.37$ and latched units saturate at $a_j=1$.

When activity is effectively constant ($\bm a_k\approx\bar{\bm a}$) and projection is inactive, the PES update reduces to the recursion for $\bm z_k=W_k\bar{\bm a}$
\begin{equation}
 \bm z_{k+1}=(1-\lambda_W)\bm z_k
 +\eta\|\bar{\bm a}\|^2[\bm y_k^*-\clip_{[-1,1]}(\bm z_k)],
 \label{eq:constantbasis}
\end{equation}
a leaky first-order filter of the teaching target, matching the idealized
contraction of Proposition~\ref{prop:pes} (Sec.~\ref{sec:properties}) with
leak and clipping restored, i.e., lagged proportional action with corner
$\pesCorner$ rad/s. Hard45’s highest instantaneous frequency is $2.1$ rad/s
\cite{sandoval2018instantaneous}, below that corner.

\emph{Conclusion.} Decoder adaptation is the dominant lever: CMG dynamics supply a bounded (graded/latched) basis on which PES shapes the correction. The leak $\lambda_W$ is the $\sigma$-modification term \cite{ioannou1996robust}, leaving the stated PES mismatch bounds.

%We distinguish decoder adaptation (PES), input-dependent conductance modulation ($\chi$), and readout gain ($G$). Toggling $\chi$ retains the base shunting term $g_0r_j^2$; it is not removal of all conductance or all CMG dynamics.
%Fixed-decoder residuals already improve some missions: in the audited direct
%X2 runs their mean reductions range from about 0.12\% to 10.27\%; the
%corresponding CF2-design range is 0.06\% to 6.80\%. Thus it is inaccurate
%to say that fixing the decoder removes \emph{all} improvement. PES produces
%the much larger 17.3--60.6\% primary reductions. Separately, toggling
%$\chi$ with the decoder mode held fixed is numerically indistinguishable
%under the tested settings. Decoder adaptation is therefore the principal lever in the tested regime.

%Readout saturation does not imply negligible conductance. In the fresh canonical X2 Full3 rerun, $r_j>0.9$ for $97.7\%$ of unit-samples and mean $c_j\tau$ is $7.08$. Yet $a_j$ clips at $r_j=10^{-5}$.
%Moreover, with $\gamma=2$ and $v_j\ge0$, $r_j\ge1/2$ implies $b_j\ge1/2$ in \eqref{eq:gate_step}; its convex update preserves that high-gate region. Conductance can therefore change membrane voltage without producing a distinguishable decoded activity.

%The resulting decoder has a useful exact reduction. On an interval with
%$\bm a_k=\bar{\bm a}$ and inactive weight projection, let
%$\bm z_k=W_k\bar{\bm a}$. Then \eqref{eq:pes} gives
%\begin{equation}
% \bm z_{k+1}=(1-\lambda_W)\bm z_k+
% \eta\|\bar{\bm a}\|^2[\bm y_k^*-\clip_{[-1,1]}(\bm z_k)].
% \label{eq:constantbasis}
%\end{equation}
%When readout clipping is inactive, this is a first-order target filter.
%For six unit activities its pole is $.995199$, corresponding to about
%$.104$ s. This is a direct-backend identity under the stated conditions,
%not equivalence of the full transient CMG or delayed Nengo dynamics.
%The recorded X2 Full3 activity is identically one after $2.02$ s; the
%reduced recursion matches its saved decoder to $8\times10^{-16}$.

%A focused direct-backend sweep uses $G\in\{1,10^3,10^5,10^6\}$ and $\chi\in\{0,1\}$ on X2 Full3 and CF2 Hardstep. Their reductions remain $56.04$--$56.97\%$ and $18.34$--$18.46\%$, respectively. In contrast, X2 Roll ranges from effectively zero at $G=1$ to $17.36\%$ at $10^5$ and $57.76\%$ at $10^6$ (also tested at $10^2,10^4$). These are fixed-$\eta$ sensitivity tests: before clipping, same-feature PES excitation scales as $\eta G^2\|\bm r\|^2$. They do not isolate neuron expressiveness. Eight selected Nengo checks agree with the direct mean errors to relative difference $1.13\times10^{-7}$.

%\begin{table}[ht]
%\centering\footnotesize
%\caption{Direct-backend controls: mean-error reduction over matched
%PD [\%]. Same $W_0$, PES teacher/update, leak, caps, and allocation;
%canonical CMG uses $G=10^5$. One deterministic run per cell.}
%\label{tab:simplebasis}
%\setlength{\tabcolsep}{4pt}
%\begin{tabular}{lrrr}
%\toprule
%Case & CMG & Paired-linear & Constant\\\midrule
%X2 Roll &17.36&0.73&57.78\\
%X2 Full3 &56.56&6.70&56.98\\
%CF2 Hardstep &18.37&1.50&18.38\\\bottomrule
%\end{tabular}
%\end{table}

% (Optional) Intermittent-disturbance figure removed to meet the 8-page ACC limit.
%\begin{figure}[t]
%\centering
%\includegraphics[width=\linewidth]{figures/07_disturbance_and_noise/hover_intermittent.png}
%\caption{X2 hover hold under four $0.15$\,N\,m disturbance bursts (shaded)
%separated by quiet intervals. Top: attitude error. Centre: residual
%magnitude, logarithmic. Bottom: decoder displacement from $W_0$. The
%residual is at numerical precision before the first burst. This means
%negligible initial correction, not zero computation or energy cost.}
%\label{fig:hover_int}
%\end{figure}

%%%%%%%%%%%%%%%%%%%%%%%%%%%%%%%%%%%
\subsection{Neuron-Family Comparison}
\label{sec:family}
%%%%%%%%%%%%%%%%%%%%%%%%%%%%%%%%%%%%%%

The Nengo benchmark spans two models, five missions, and eleven controllers (PD plus fixed/PES variants of CMG, LIF, ALIF, Izhikevich \cite{izhikevich2003}, and LIFRate). All share the same bounded residual and actuator interface; only the neuron dynamics and native-state-to-activity map differ.
Across the full benchmark matrix (figure omitted), two patterns hold. First, PES is the
dominant lever in every family: fixed decoders rarely match their PES
counterparts, even with more expressive dynamics, extending
Sec.~\ref{sec:ablation} to every basis. Second, lower error is not free:
some Izhikevich--PES cells reach the smallest error, but at an average
normalized residual RMS near 0.37 against 0.12 for CMG--PES, mission by mission in
Table~\ref{tab:all}, and with larger decoder movement, so equal output caps do not
imply equal realized effort.
Proposition~\ref{prop:pes} explains the flat rows: contraction scales with $1-\eta\|\bm a\|_2^2$, so a basis that is silent stalls the update. LIF, ALIF and LIFRate units are active in $0.1\%$ of unit-samples on Full3 and Hard45 and under $0.05\%$ on Roll and Yaw360, whereas CMG (resting at $a_j\approx0.37$ on its logistic tail) and Izhikevich are active in ${\ge}99.9\%$. Persistent activity is therefore necessary for effective PES adaptation in the tested regime, but it does not by itself set effort, since Izhikevich shares it and still spends three times the authority. CMG--PES thus attains comparable tracking at markedly lower realized effort and
more conservative decoder motion.

%Across neuron models, improved tracking typically requires greater residual effort; CMG--PES lies near a favorable point on this effort--accuracy tradeoff.
%The Nengo benchmark contains two models, five missions, and eleven controllers: PD and fixed/PES variants of CMG, LIF, ALIF, Izhikevich \cite{izhikevich2003}, and LIFRate. All share the same bounded residual, saturation, and allocation interface; only the neuron dynamics and the mapping from native state to activity differ.

%Figure~\ref{fig:mech_dashboard} summarizes the full matrix on both
%platforms: mean tracking error and realized residual effort for every
%family and learning mode. Two consistent patterns emerge. First, PES adaptation is the dominant lever across families: fixed
%decoders rarely match their PES counterparts, even when the underlying
%neuron dynamics are more expressive. Second, lower error does not come for
%free. Some Izhikevich--PES cells achieve the smallest mean error, but the
%audit reports a substantially larger average normalized residual RMS
%(approximately 0.37 versus 0.12 for CMG--PES) and larger decoder movement,
%so equal output caps do not imply equal realized correction effort.

%Sparse LIF/ALIF activity in some cells limits the effective feature basis at
%the tested operating point, which can blunt online learning without
%implying that those families cannot learn under different encodings or gain
%settings. The supported conclusion is therefore a tracking--effort tradeoff:
%CMG--PES attains comparable tracking at markedly lower realized residual
%effort and more conservative decoder motion in this benchmark.

%In Fig.~\ref{fig:full3_neurons}, gate/activity saturation and rotor
%clipping are distinct phenomena. Labels referring to a residual torque
%denote the requested difference $\bm\delta$, not its post-allocation effect.

% [v10] One two-column figure replaces the two single-column dashboards.
% The 181 mm source was previously scaled into an 88 mm column, a 2.05x
% reduction that left its labels near 3 pt. Drawn at 1:1 here.

\begin{comment}
\begin{figure*}[t]
    \centering
    \includegraphics[width=\textwidth]{figures/08_Nengo_neuron_family/mechanism_both_platforms.pdf}
    \caption{Mechanism separation across the Nengo neuron-family benchmark,
    X2 (top) and CF2 (bottom). Left: mean $q$-error improvement over the PD
    baseline. Right: normalized residual effort, the RMS residual as a
    fraction of the per-axis limit. Rows are neuron family and learning
    mode, columns are missions; every cell uses the same bounded residual,
    saturation, and allocation interface. Two patterns hold on both
    platforms: PES adaptation, not the neuron model, separates the improving
    rows from the flat ones; and the families that match or exceed
    CMG--PES on error do so at three to twenty times its residual effort.}
    \label{fig:mech_dashboard}
\end{figure*}
\end{comment}

%%%%%%%%%%%%%%%%%%%%%%%%%%%%%%%%%%%%%%%%%%%%%%%%%%%%%%%%
\subsection{Measurement Noise and Disturbance Rejection}
\label{sec:noise}

Each case was rerun at two synthetic noise levels with two paired seeds. At each step,
independent zero-mean Gaussian samples form a small Euler-angle quaternion
left-multiplied into $q$, and additive samples corrupt each rate component. The
moderate level is $0.010$\,rad and $0.025$\,rad\,s$^{-1}$, the realistic level
$0.020$\,rad and $0.050$\,rad\,s$^{-1}$; neither is a calibrated sensor specification.
We evaluate $\theta_q$ from the \emph{corrupted} quaternion used by the controller, so
the noise rows of Table~\ref{tab:all} measure noisy-feedback error and inherit its floor.

All noisy reductions stay positive, spanning $18.08$--$55.75\%$ at the higher level.
Yaw360 loses about $12.8$ percentage points against clean feedback (closest to the
noise floor), whereas Roll's reduction increases. Seed standard deviations are
$.001$--$.101$ pp (moderate) and $.002$--$.031$ pp (realistic): limited repeatability,
not a confidence interval.

\begin{comment}
\begin{figure}[t]
\centering
\includegraphics[width=\linewidth]{figures/02_neural_activity/cf2-rotor_hardstep_neurons.png}
\caption{CF2 Hardstep: internal CMG signals on the same trajectory as
Fig.~\ref{fig:states}. The plot separates membrane state, raw gate $r_j$, and
the clipped activity $a_j=\clip_{[0,1]}(Gr_j)$ presented to the decoder,
illustrating the readout saturation of Sec.~IV-D under the one mission that
also drives the rotors to a bound.}
\label{fig:cf2_hardstep_neurons}
\end{figure}
\end{comment}

An equilibrium disturbance-rejection test applies
$\bm\tau_{\rm d}=d_b\varrho(t)[\sin(1.31t+.2),\allowbreak
.8\sin(1.73t+1.1),\allowbreak.6\sin(2.17t+2)]\T$ in body axes, $\varrho$ alternating
5-s on/off intervals after 5 s of quiet with 0.25-s raised-cosine ramps and
$d_b=.15$ N\,m on X2 (the first component's amplitude). The four-burst 45-s run
cuts burst-phase mean error from $11.50^\circ$ to $6.65^\circ$ ($42.2\%$) and
post-burst quiet error from $1.21^\circ$ to
$.99^\circ$ ($18.0\%$); the CF2 rerun ($d_b=.00015$ N\,m, a selected scaling)
falls from $.141199$ to $.070259$ rad ($50.24\%$).

%The 180-s X2 record contains eighteen bursts and seventeen complete post-burst quiet intervals, not sixteen bursts. Decoder displacement reaches $.797$ at the final sample, with maximum absolute weight $1.089$.
%One quiet interval is about $4\%$ worse than PD. Continued weight change and nonzero quiet-phase errors do not establish convergence. Because learning and leakage remain active, and the multi-sine phase is not reset between bursts, this study shows persistent adaptation, not causal proof of reuse without relearning. A reset-decoder or frozen-decoder control is needed to test that hypothesis.

%%%%%%%%%%%%%%%%%%%%%%%%%%%%%%%%%%%%%%%%%%%%%%%%
\subsection{Failure Regime: Abrupt Hardstep Commands}
\label{sec:ladder}

Hardstep is kinematically dominated by slew transients (piecewise-constant
reference gives $\bm\omega_d\equiv\bm0$), so mean-error improvements are
muted even when settling improves; against fixed PD, rotor commands reach a
bound only on Hardstep ($S=.016/.239$ on X2 and $.137/.264$ on CF2,
PD/CMG). At $\rho=3$ Hardstep is the only cell where the residual degrades:
the latched-decoder dynamics \eqref{eq:constantbasis} act as a first-order
target filter (Sec.~\ref{sec:properties}), so the residual lags a
high-bandwidth loop and, under torque-box saturation, can only hurt when it
opposes $\bm u_0$.

A simple remedy is to gate the residual on saturated, opposed axes:
$s_i\leftarrow0$ when $|u_{0,i}|=\ell_i$ and $\operatorname{sgn}s_i\neq\operatorname{sgn}u_{0,i}$. This preserves $|\bm u-\bm u_0|\le\beta\bm\ell$ and exact baseline recovery (Proposition~\ref{prop:interface}); we leave a systematic evaluation under noise/disturbance/mismatch to future work (all results here use the ungated law).

\begin{comment}
\begin{figure}[t]
\centering
\includegraphics[width=\linewidth]{figures/02_neural_activity/attitude_detail/cf2_rotor_hardstep_attitude.png}
\caption{CF2 Hardstep mission: PD baseline versus CMG--PES augmentation on
the abrupt step reference. This is the hardest case in the matrix, the
smallest reduction ($18.4\%$), and the only one whose rotor commands reach a
bound ($S=.137/.264$). The residual stays inside its
authority envelope and the augmented loop remains better than the baseline
throughout.}
\label{fig:states}
\end{figure}
\end{comment}

\section{Conclusions and Future Work}
\label{sec:conclusion}
% [v7] Rewritten: states the positive result once and plainly, keeps the
% limits in one place instead of after every sentence, and separates the
% analog outlook from the claims.
% [v12] Result paragraph now cites the tuned PD-feedforward comparison
% (rho=3) alongside fixed PD, matching Sec. IV-B's lead comparison.

Six continuous CMG units with an 18-parameter PES-adapted decoder form a bounded residual that augments a fixed quaternion PD baseline without changing gains or the actuator interface. Over ten matched 45-s MuJoCo cases on two platforms, mean quaternion-angle error drops by $17.3$--$60.6\%$ against fixed PD and by $6.5$--$35.6\%$ (geometric mean) against a tuned PD-feedforward baseline at $\rho=3$, with exact baseline recovery at zero authority; direct and Nengo realizations match to $2.013\times10^{-6}$ relative mean error.

Mechanism: PES (not $\chi$) drives the improvement; gates latch and saturate activity, while sub-threshold units provide a graded tail. Across five neuron families PES is the consistent differentiator and neuron dynamics mainly trade tracking for residual effort; $G$ (resting activity $\approx0.37$) is a key sensitivity.

Limitations: simulation study only (noise sensitivity bounded in Sec.~\ref{sec:noise}); model mismatch/hardware effects are not addressed; retention is limited; bounded commands/states/weights are implementation properties, not stability proofs. The bounded-residual interface \eqref{eq:residual} is benchmarked only against other learned neural bases (Table~\ref{tab:all}); it is not compared against a non-neuromorphic adaptive law (e.g., RLS or normalized-LMS) under the identical bounded interface, so the reported gains isolate the CMG basis among neural alternatives and do not by themselves show an advantage over adaptive control generally.

Future work: stronger/matched baselines (incl. integral action), evaluation of saturation-gated residuals (Sec.~\ref{sec:ladder}) under noise/disturbance/mismatch, matched-effort comparisons to non-neural bounded bases, sweeps of $G$ and $\chi$, longer retention and held-out tests with frozen weights, 6-DOF extension, and hardware validation; the wrapper extends to nonlinear/optimal baselines with hard actuator constraints.

%A bounded neuromorphic residual can augment a fixed quaternion PD baseline without changing its gains or actuator interface. Using six continuous CMG units with eighteen PES-adapted decoder weights, the augmented loop reduces mean quaternion-angle tracking error by $17.3$--$60.6\%$ over ten matched 45-s MuJoCo cases on two modeled platforms, while respecting rotor saturation and recovering the baseline exactly at zero residual authority. Direct and Nengo realizations agree to $2.013\times10^{-6}$ relative mean error.

%Mechanism studies show that, in the tested regime, decoder adaptation is the dominant source of improvement: fixed-decoder residuals yield only modest reductions, and toggling the input-dependent conductance term ($\chi$) with the decoder mode held fixed is numerically indistinguishable. Across CMG, LIF/ALIF, Izhikevich, and LIFRate bases, PES is the consistent differentiator, while neuron dynamics mainly set an effort--accuracy tradeoff (matching error often requires substantially more residual effort). The observed readout saturation $a_j=\clip(Gr_j,0,1)$ suggests $G$ is a key sensitivity parameter.

%Limitations: results are in simulation (Sec.~\ref{sec:noise} bounds noise sensitivity, but model mismatch and hardware effects are not addressed); the repeated-disturbance retention horizon is limited; and bounded commands, states, and weights are implementation properties rather than closed-loop stability proofs.

%Future work will include integral-action baselines, matched-effort comparisons against non-neural bounded bases, systematic sweeps of $G$ and $\chi$, longer-horizon retention studies, held-out evaluations with frozen learned weights, extension to full 6-DOF systems, and, most importantly, validation on physical robotic platforms, including aerial, ground, and other dynamical systems. The same bounded-residual wrapper can also be extended to nonlinear and optimal-control baselines subject to hard actuator constraints.

% Finalize human authors, affiliations, funding, and any required tool-use
% disclosure before submission. No invented author identities are supplied.
\bibliographystyle{IEEEtran}
\bibliography{references_acc2027}
\end{document}
